\documentclass[11pt,letterpaper]{article}
\usepackage{fontspec}
\usepackage[margin=1in]{geometry}
\usepackage{amsmath,amssymb,amsthm,mathtools}
\usepackage{unicode-math}
\setmainfont{Latin Modern Roman}
\setmonofont{Latin Modern Mono}
\setmathfont{Latin Modern Math}
\usepackage{microtype,booktabs,array,enumitem}
\usepackage[hyphens]{url}
\usepackage[hidelinks]{hyperref}
\usepackage{fancyhdr}
\pagestyle{fancy}\fancyhf{}
\fancyhead[L]{\small Prefix stabilizer maxima}\fancyhead[R]{\small Manuscript version 3}
\fancyfoot[C]{\thepage}\renewcommand{\headrulewidth}{0pt}
\lefthyphenmin=3\righthyphenmin=3
\setlength{\headheight}{14pt}\setlength{\emergencystretch}{2em}
\newtheorem{theorem}{Theorem}
\newtheorem{lemma}[theorem]{Lemma}
\newtheorem{proposition}[theorem]{Proposition}
\theoremstyle{definition}\newtheorem{definition}{Definition}
\newcommand{\Stab}{\operatorname{Stab}}
\newcommand{\Max}{\operatorname{Max}}
\newcommand{\ind}{\mathbf{1}}
\newcommand{\D}{\mathcal D}
\newcommand{\C}{\mathcal C}
\newcommand{\hash}[1]{{\small\ttfamily\nolinkurl{#1}}}
\setlist[enumerate]{itemsep=3pt,topsep=5pt}
\title{\bfseries Prefix stabilizer maxima for two Sylow factors}
\author{Research manuscript\quad Version 3 candidate}
\date{11 October 2026}
\begin{document}\maketitle\thispagestyle{plain}
\begin{abstract}
For a labeled ordering of a finite group, sum the orders of the left setwise stabilizers of all nonempty prefixes, subtracting one at each prefix. Let $H$ be a finite $q$-group of order $q^a$ and $P$ a finite $p$-group of order $p^b$, where $p\ne q$. The maximum of this score on $H\times P$ is
\[
q^ap^b\left[a\left(1-\frac1q\right)+b\left(1-\frac1p\right)\right].
\]
Both factors may be nonabelian. An endpoint closure and core argument supplies strict clipping bounds for a numerical hierarchy, even when the checkpoint stabilizers have incomparable $P$-projections. Whole-gap densification and an exact rescaling identity prove the upper bound; consecutive maximal-subgroup coset recursion attains it. For $P=C_{p^b}$ and, separately, for $P=C_p^2$, the maximizers are exactly these recursive orders, the proper-prefix stabilizers recover their top subgroup uniquely, and labeled counting recurrences follow. The corresponding sufficient escape deficits are $p^{b-1}(p-1)$ and $p-1$. A $D_8\times C_3^2$ example explains why the cyclic escape argument cannot be transferred unchanged. General $P$ optimizer classification, sharp escape values, and historical priority remain unresolved here.
\end{abstract}

\paragraph{Version and review status.}
This version integrates the independently accepted noncyclic maximum theorem and elementary-abelian rank-two structural theorem into the cyclic-prime-power manuscript. The mathematical source has independent project acceptance, not external peer review or formal certification. \textbf{Version 3 is a candidate manuscript pending independent document-fidelity review.} The eleven-page version 1 and fourteen-page version 2 remain unchanged. Appendix~\ref{app:provenance} and the accompanying package bind the exact accepted sources and preserved versions.

\section{Statement and conventions}
Let $q,p$ be distinct primes, let $|H|=n=q^a$ with $H$ a finite $q$-group, and let $P$ be any finite $p$-group of order $p^b$, where $a,b\ge0$. Put $G=H\times P$ and $N=np^b$. For the cyclic specialization write $P_b=C_{p^b}$, with $P_0=1$. For an ordering $\mathbf{x}=(x_1,\ldots,x_N)$ of the labeled elements, define
\begin{align}
T_j&=\{x_1,\ldots,x_j\},& S_j&=\{g\in G:gT_j=T_j\},\label{eq:prefix}\\
Z_G(\mathbf{x})&=\sum_{j=1}^{N}(|S_j|-1),& B(\mathbf{x})&=\langle S_1,\ldots,S_{N-1}\rangle.\label{eq:score}
\end{align}
The action is by left multiplication. The full prefix is included in $Z_G$ and excluded from $B$. Write $M(K)$ for the maximum score on a finite group $K$ and $A(K)$ for the number of labeled maximizing orders. No quotient by translations or automorphisms is taken.

\begin{definition}
A recursive maximal-subgroup right-coset ordering of a nontrivial group chooses a maximal proper subgroup $J$, places its distinct right cosets consecutively, and recurses independently inside each coset. The singleton has its unique ordering. The identification $Jx\to J$, $hx\mapsto h$, defines the local recursion. Changing representatives merely changes these coordinates by right translation, which preserves left stabilizers.
\end{definition}

Write
\begin{equation}
V_{a,b}=np^b\left[a\left(1-\frac1q\right)+b\left(1-\frac1p\right)\right].\label{eq:value}
\end{equation}
The three theorem scopes below are deliberately separate.

\begin{theorem}[Two arbitrary Sylow factors]\label{thm:two-sylow}
For all $a,b\ge0$ and all finite groups $H,P$ of orders $q^a,p^b$, respectively, where $p\ne q$,
\[
M(H\times P)=V_{a,b}.
\]
Recursive maximal-subgroup right-coset orders attain this value. This theorem alone does not classify all maximizing orders, recover a unique top subgroup, give their total count, or give an escape deficit for arbitrary $P$.
\end{theorem}

\begin{theorem}[Cyclic second factor]\label{thm:main}
For $G=H\times P_b$, where $P_b=C_{p^b}$ and $a,b\ge0$, the following additional assertions hold.
\begin{enumerate}
\item An ordering maximizes $Z_G$ if and only if it is a recursive maximal-subgroup right-coset ordering.
\item For nontrivial $G$, the top subgroup of every maximizing ordering is unique and equals $B(\mathbf{x})$.
\item With $A(1)=1$, for $b\ge1$,
\begin{equation}
A(H\times P_b)=p!\,A(H\times P_{b-1})^p
 +q!\sum_{M\in\Max(H)}A(M\times P_b)^q.\label{eq:count}
\end{equation}
The sum is empty when $H=1$. For $b=0$ and $H\ne1$,
\begin{equation}
A(H)=q!\sum_{M\in\Max(H)}A(M)^q.\label{eq:qcount}
\end{equation}
\item If $a,b\ge1$ and $B(\mathbf{x})=G$, then
\begin{equation}
Z_G(\mathbf{x})\le V_{a,b}-p^{b-1}(p-1).\label{eq:escape}
\end{equation}
\end{enumerate}
\end{theorem}

\begin{theorem}[Elementary abelian second factor of rank two]\label{thm:ranktwo}
Let $P=C_p^2=C_p\times C_p$ and $a\ge0$. The maximizing orders on $H\times P$ are exactly the recursive maximal-subgroup right-coset orders. Every nontrivial maximizing order has a unique top subgroup, equal to $B(\mathbf{x})$. With $A(1)=1$,
\begin{equation}
A(H\times C_p^2)=(p+1)p!\,A(H\times C_p)^p
+q!\sum_{M\in\Max(H)}A(M\times C_p^2)^q.\label{eq:rankcount}
\end{equation}
The sum is empty when $H=1$. If $a\ge1$ and $B(\mathbf{x})=H\times C_p^2$, then
\begin{equation}
Z_G(\mathbf{x})\le V_{a,2}-(p-1).\label{eq:rankescape}
\end{equation}
\end{theorem}

Both structural statements allow independent optimal local choices in distinct cosets and terminate at singletons. Counts are labeled counts, with no translation or automorphism quotient. The singleton has value $0$, one order, and no top subgroup. The escape estimates are sufficient deficits, not sharp escaping maxima, and neither is asserted on a prime-power axis. The results allow $p=2<q$ and arbitrary nonabelian $H$. Theorem~\ref{thm:two-sylow} also allows nonabelian $P$. No theorem here treats three distinct prime divisors or all finite nilpotent groups. The separate stronger historical $D_8\times C_p$ escape estimate is retained in Section~\ref{sec:d8}.

\section{Threshold checkpoints and endpoint closure}\label{sec:checkpoints}
Every subgroup $S\le H\times P$ splits as $L\times Q$. Indeed, coprime-order powers of any $(h,z)\in S$ isolate $(h,1)$ and $(1,z)$. Applying this to every element gives the decomposition, without commutativity within either factor. Write
\[
S_j=L_j\times Q_j,\qquad \beta_j=|L_j|,\qquad |Q_j|=p^{r(j)}.
\]
Define the threshold checkpoint sets
\[
C_s=\{0\}\cup\{j:1\le j\le N,\ |Q_j|\ge p^s\}.
\]
They use all actual checkpoints and satisfy
\begin{equation}
C_b\subseteq C_{b-1}\subseteq\cdots\subseteq C_1\subseteq C_0=\{0,\ldots,N\},
\quad C_s\subseteq p^s\mathbb Z,\quad 0,N\in C_s.\label{eq:hierarchy}
\end{equation}
The divisibility follows because $S_j$ acts freely on $T_j$, so $|S_j|\mid j$. These are nested sets of indices, not a common chain of actual $P$-subgroups. Distinct checkpoints may have incomparable projections. Every finite $p$-group has subgroups of each order dividing its order: recursively take normal prime-index maximal subgroups, as proved in Lemma~\ref{lem:normal}. Consequently a level-$s$ checkpoint is invariant under some order-$p^s$ subgroup of $P$. When $P$ is cyclic, this specializes to its unique subgroup of that order.

For $\D=\{q,q^2,\ldots,q^a\}$ and $c_d=d-d/q$, telescoping gives
\begin{equation}
\beta-1=\sum_{\substack{d\in\D\\d\mid\beta}}c_d,
\qquad \gcd(t,n)=1+\sum_{\substack{d\in\D\\d\mid t}}c_d.\label{eq:weights}
\end{equation}
The set $\D$ is empty when $n=1$.

\begin{lemma}[Endpoint closure and core]\label{lem:closure}
Fix $j$ with $r=r(j)<s\le b$, and let $a_s<j<z_s$ be its consecutive enclosing $C_s$ indices. If $d\in\D$ divides $\beta_j$, then
\begin{equation}
d\mid j/p^r,\qquad
p^sd\left\lceil\frac{a_s}{p^sd}\right\rceil<j
 <p^sd\left\lfloor\frac{z_s}{p^sd}\right\rfloor.\label{eq:clips}
\end{equation}
This holds simultaneously at every enclosing level, without requiring the endpoint $P$-subgroups to agree or to be comparable.
\end{lemma}
\begin{proof}
Choose $A_s,D_s\le P$ of order $p^s$ preserving $T_{a_s},T_{z_s}$, respectively. At an empty or full endpoint any such subgroup is available. Put $L=L_j$ and
\[
F=\bigcup_{\ell\in L}\ell T_{a_s},\qquad
E=\bigcap_{\ell\in L}\ell T_{z_s}.
\]
The $L$-invariance of $T_j$ gives
\[
T_{a_s}\subseteq F\subseteq T_j\subseteq E\subseteq T_{z_s}.
\]
Left multiplication by $L$ permutes the union and intersection factors. The direct factors commute, so $A_s$ preserves each $\ell T_{a_s}$ and $D_s$ preserves each $\ell T_{z_s}$. Thus $F$ is invariant under $L\times A_s$, while $E$ is invariant under $L\times D_s$. Their free translation actions show that both cardinalities are multiples of $\beta_jp^s$, and hence of $dp^s$. Therefore
\[
p^sd\left\lceil\frac{a_s}{p^sd}\right\rceil
\le |F|\le j\le |E|
\le p^sd\left\lfloor\frac{z_s}{p^sd}\right\rfloor.
\]
Equality at the outer left bound forces $F=T_j$, making $T_j$ invariant under $A_s$ and contradicting $r<s$. Equality at the right similarly forces $E=T_j$ and $D_s$-invariance. Finally $\beta_jp^r\mid j$ gives $d\mid j/p^r$.
\end{proof}

At a top checkpoint $j=p^bk$, the prefix is $A\times P$ for a subset $A\subseteq H$ of size $k$. Its $H$-stabilizer acts freely on $A$, so $\beta_j\mid\gcd(k,n)$.

\section{The all enclosing numerical invariant}\label{sec:clips}
Now allow \emph{any numerical hierarchy} $\C=(C_0,\ldots,C_b)$ satisfying~\eqref{eq:hierarchy}, whether realizable or not. For $j>0$, define $r(j)=\max\{s:j\in C_s\}$ and put $r=r(j)$. Then define
\begin{equation}
h_{\C}(j)=
\begin{cases}
\gcd(j/p^b,n),&r=b,\\
1+\displaystyle\sum_{d\in\D}c_d I_{\C}(d,j),&r<b,
\end{cases}\label{eq:cap}
\end{equation}
where $I_{\C}(d,j)=1$ exactly when $d\mid j/p^r$ and all of~\eqref{eq:clips} hold, for every $s>r$. Otherwise it is zero. Put
\begin{equation}
U_b(\C)=\sum_{j=1}^{N}\left[p^{r(j)}h_{\C}(j)-1\right].\label{eq:invariant}
\end{equation}
The cap need not be realizable or a subgroup order. For the actual checkpoint hierarchy,~\eqref{eq:weights}--\eqref{eq:clips} imply $\beta_j\le h_{\C}(j)$. Retain the entire top-endpoint deficit
\begin{equation}
D_{\rm top}=\sum_{\substack{k>0\\p^bk\in C_b}}
 p^b\left[\gcd(k,n)-\beta_{p^bk}\right].\label{eq:deficit}
\end{equation}
Its summands are nonnegative and its final summand is zero. Thus the strengthened upper bound is
\begin{equation}
Z_G(\mathbf{x})\le U_b(\C)-D_{\rm top}.\label{eq:strongcap}
\end{equation}
The induction will preserve the numerical hierarchy, every enclosing clip, and this unchanged deficit from the actual top checkpoints.

\section{Whole gap densification}\label{sec:dense}
Assume $b\ge1$. Let $pu<pv$ be consecutive $C_1$ indices and write $R=v-u\ge1$. There is no positive-level checkpoint in their open physical interval. Therefore for every $s\ge2$, its enclosing $C_s$ endpoints $a_s,z_s$ are constant throughout this interval. For $d\in\D$ put
\begin{align}
L_d&=\max\left(\{u\}\cup\left\{p^{s-1}d\left\lceil\frac{a_s}{p^sd}\right\rceil:2\le s\le b\right\}\right),\label{eq:L}\\
U_d&=\min\left(\{v\}\cup\left\{p^{s-1}d\left\lfloor\frac{z_s}{p^sd}\right\rfloor:2\le s\le b\right\}\right).\label{eq:U}
\end{align}
When $b=1$, these are just $u,v$. All higher-level clipping endpoints in these expressions are multiples of $d$. If $L_d\ge U_d$, define $N_d=R_d=0$. Otherwise set
\begin{equation}
N_d=\max\left(0,p\left[\left\lfloor\frac{U_d}{d}\right\rfloor-
\left\lceil\frac{L_d}{d}\right\rceil\right]-1\right),
\qquad R_d=\#\{t\in\mathbb Z:L_d<t<U_d,\ d\mid t\}.\label{eq:NR}
\end{equation}
The old open interval contributes exactly $\sum_d c_dN_d$ to $U_b$. Its indices all have level zero. Eligibility for $d$ requires divisibility by $d$ and the strict level-one clip
\[
pd\lceil u/d\rceil<j<pd\lfloor v/d\rfloor.
\]
Intersecting with all higher clips gives
\[
pd\lceil L_d/d\rceil<j<pd\lfloor U_d/d\rfloor,
\]
whose eligible multiples of $d$ are counted by $N_d$.

Densify $C_1$ by inserting \emph{all} missing inner ranks $pt$, $u<t<v$, at once, while leaving every $C_s$, $s\ge2$, unchanged. The inserted indices have level one and their eligibility for $d$ is exactly $d\mid t$ and $L_d<t<U_d$. Every remaining level-zero index has no $q$-weight in its cap: it lies between adjacent new $C_1$ ranks, too close to enclose the required distinct $d$-multiples for $d\ge q\ge2$. Its cap $h_{\C}$ is therefore $1$ and its score contribution is zero. The new open-interval contribution is
\[
(p-1)(R-1)+p\sum_d c_dR_d.
\]
The new minus old difference is
\begin{equation}
\Delta=(p-1)(R-1)+\sum_d c_d(pR_d-N_d).\label{eq:delta}
\end{equation}
Here is a complete arithmetic check of its sign. For $L_d\ge U_d$, the summand is zero. Otherwise there are three cases:
\begin{enumerate}
\item Both endpoints are divisible by $d$. If $w=(U_d-L_d)/d\ge1$, then $R_d=w-1$, $N_d=pw-1$, and $pR_d-N_d=-(p-1)$.
\item Exactly one endpoint is divisible by $d$. Then $N_d=\max(0,pR_d-1)$; the difference is $0$ for $R_d=0$, and $1$ otherwise.
\item Neither endpoint is divisible by $d$. The difference is $0$ for $R_d=0$, $p$ for $R_d=1$, and $p+1$ for $R_d\ge2$.
\end{enumerate}
Only the first case is negative, and it requires $d\le U_d-L_d\le R$. Put
\[
d_* = \max\{q^i:0\le i\le a,\ q^i\le R\}.
\]
In particular $d_*\le\min(R,n)$, even for gaps with $R>n$. The telescoping identity gives
\[
\sum_{\substack{d\in\D\\d\le R}}c_d=d_*-1\le R-1,
\]
and hence
\begin{equation}
\Delta\ge(p-1)\bigl[R-1-(d_*-1)\bigr]=(p-1)(R-d_*)\ge0.\label{eq:deltasign}
\end{equation}
All pre-existing checkpoint caps stay unchanged, since a positive-level cap only uses levels strictly above its own. Repeating over the gaps, let $\C^{\rm dense}$ have $C_1$ equal to all multiples of $p$. Then
\begin{equation}
U_b(\C)\le U_b(\C^{\rm dense}).\label{eq:densification}
\end{equation}
This compares whole gaps of numerical bounds. It assumes neither single-checkpoint insertion monotonicity nor realizability of the densified hierarchy by a group ordering.

\section{Exact quotient identity at arbitrary depth}\label{sec:quotient}
In $\C^{\rm dense}$, every nonmultiple of $p$ has zero score contribution. At each remaining index $j=pt$, decrement its level by one and divide higher checkpoint indices by $p$. The resulting hierarchy for total length $np^{b-1}$ is
\begin{equation}
C'_0=\{0,\ldots,np^{b-1}\},\qquad
C'_s=\{j/p:j\in C_{s+1}\}\quad(1\le s\le b-1).\label{eq:quotient}
\end{equation}
The clipping, divisibility, and top-$H$ caps correspond exactly. The retained $q$-power data remain $\D=\{q,\ldots,q^a\}$. At each positive quotient index,
\[
p^rh-1=p\bigl[p^{r-1}h-1\bigr]+(p-1),
\]
so there is an exact identity
\begin{equation}
U_b(\C^{\rm dense})=pU_{b-1}(\C')+(p-1)np^{b-1}.\label{eq:quotientidentity}
\end{equation}
For depth zero, the definition is $U_0=\sum_{k=1}^{n}[\gcd(k,n)-1]$. Equation~\eqref{eq:quotientidentity} is an identity of the numerical invariant; it does not invoke a group theorem on an artificial quotient order.

The intermediate length $np^{b-1}$ may have both prime divisors. It is never substituted for $n$ in $\gcd(t,n)$ or in $\D$. The $H$ rank base remains the prime power $n=q^a$; only the numerical $p$-depth decreases.

Apply~\eqref{eq:densification} and~\eqref{eq:quotientidentity} $b-1$ times. Every additive term, after preceding scalings, is $(p-1)np^{b-1}$. Therefore
\begin{equation}
U_b(\C)\le p^{b-1}U_1(\C^{\rm top})+(b-1)(p-1)np^{b-1},\label{eq:collapse}
\end{equation}
where $C^{\rm top}_1=\{j/p^{b-1}:j\in C_b\}$. The actual top checkpoints never change. The actual deficit $D_{\rm top}$ in~\eqref{eq:strongcap} remains subtractable after every replacement; no endpoint deficit is discarded and no earlier stabilizer is assumed to persist at later prefixes.

The same elimination restricts to a consecutive top interval $p^bk<p^bl$. Its additive term is
\[
(b-1)(p-1)p^{b-1}(l-k).
\]
This interval form will give the exact outer edge loss.

\section{The outer q adic edge and upper bound}\label{sec:edge}
Put $h(k)=\gcd(k,n)$, $s(k)=\sum_{t=1}^{k}h(t)$, and $s(0)=0$. For arbitrary integers $0\le k<l$ define
\begin{align}
w_d(k,l)&=\max\left(0,\left\lfloor\frac ld\right\rfloor-\left\lceil\frac kd\right\rceil\right),\label{eq:w}\\
E_p(k,l)&=p\left[h(l)+\sum_{d\in\D}c_dw_d(k,l)\right]
-\left[1+\sum_{\substack{d\in\D\\w_d(k,l)>0}}c_d\right],\label{eq:edge}\\
F_p(k)&=ps(k)-k.\label{eq:potential}
\end{align}
At one level, the strict clips count $\max(0,pw_d-1)$ eligible open-interval indices for $d$, and the complete right endpoint contributes $ph(l)-1$. Thus $U_1$ is exactly the sum of $E_p$ over consecutive top ranks. This includes composite prime powers $n=q^a$ and the empty $\D$ case.

For completeness, the one-level calculation is reproduced. Put $m=l-k$, $g=\gcd(k,l,n)$, and let $X$ consist of those $d\in\D$ for which neither endpoint is divisible by $d$ but an interior $d$-multiple exists. Then
\begin{align}
F_p(l)-F_p(k)-E_p(k,l)&=p\alpha+\gamma,\label{eq:loss}\\
\alpha&=m-g+\sum_{d\in X}c_d,\label{eq:alpha}\\
\gamma&=1-m+\sum_{\substack{d\in\D\\w_d>0}}c_d.\label{eq:gamma}
\end{align}
Indeed, if $N_d$ counts strict interior $d$-multiples, expanding $h(t)$ gives
\[
\alpha=m-1+\sum_d c_d(N_d-w_d).
\]
The difference $N_d-w_d$ is $-1$ for two divisible endpoints; it is $+1$ for two indivisible endpoints with an interior multiple; it is zero otherwise. The negative weights sum to $g-1$, proving~\eqref{eq:alpha}.

Since $g\mid m$, if $m>g$ then $m\ge2g$, and
\begin{align}
p\alpha+\gamma&\ge p(m-g)+1-m\notag\\
&=(p-1)(m-2g)+(p-2)g+1\ge p-1.\label{eq:positive}
\end{align}
If $m=g$, the interval is aligned $q$-adic: its endpoints are consecutive multiples of the $q$-power $m$. For $d\le m$, both endpoints are divisible and $w_d=m/d$; for $d>m$, there is no interior $d$-multiple. Consequently $X$ is empty, $\alpha=\gamma=0$, and
\begin{equation}
F_p(l)-F_p(k)-E_p(k,l)\ge0,
\quad\text{with equality exactly when }l-k=\gcd(k,l,n).\label{eq:zero}
\end{equation}
Every nonzero loss is at least $p-1$, including $p=2$. No part of this calculation uses $l\le n$; the fixed prime-power base remains $n$. This extended-rank form will be needed for Theorem~\ref{thm:ranktwo}.

For $b\ge1$ define
\begin{align}
E_b(k,l)&=p^{b-1}E_p(k,l)+(b-1)p^{b-1}(p-1)(l-k),\label{eq:collapsededge}\\
F_b(k)&=p^{b-1}F_p(k)+(b-1)p^{b-1}(p-1)k\notag\\
&=p^bs(k)+\bigl[bp^{b-1}(p-1)-p^b\bigr]k.\label{eq:collapsedpotential}
\end{align}
Equation~\eqref{eq:collapse} and its interval form give
\begin{equation}
Z_G\le\sum_{\text{consecutive top ranks }k,l}E_b(k,l)-D_{\rm top}.\label{eq:outerbound}
\end{equation}
The exact loss is
\begin{equation}
F_b(l)-F_b(k)-E_b(k,l)=p^{b-1}\bigl[F_p(l)-F_p(k)-E_p(k,l)\bigr].\label{eq:scaledloss}
\end{equation}
Thus it has the same zero criterion, and every nonzero loss is at least $p^{b-1}(p-1)$. No unproved propagation of losses through new layers is needed.

Counting the $q$-adic valuations of $1,\ldots,n$ gives
\[
s(n)=n+\sum_{i=1}^{a}(q^i-q^{i-1})\frac{n}{q^i}
=n+an\left(1-\frac1q\right).
\]
The outer potential telescopes from $0$ to $n$, and hence
\begin{equation}
Z_G\le F_b(n)=V_{a,b}.\label{eq:upper}
\end{equation}
This proves $U_b(\C)\le V_{a,b}$ for every admissible numerical hierarchy, whether group-realizable or not, and establishes the upper bound for arbitrary $P$ and every $b\ge1$. The $b=0$ group case is proved independently in Section~\ref{sec:bases}.

\section{Strict escape for a cyclic second factor}\label{sec:escape}
In this section take $P=P_b=C_{p^b}$ and identify $P_r$ with its unique subgroup of order $p^r$. Assume $a,b\ge1$ and $B(\mathbf{x})=G$. There must be a proper top checkpoint. Otherwise every proper-prefix stabilizer has $p$-part contained in the unique $P_{b-1}$, and the subgroup they generate lies in $H\times P_{b-1}<G$. This step depends on the cyclic factor: several smaller projections cannot together generate $P_b$.

Set $r_0=n/q$. A top-rank path consisting entirely of zero-loss edges either is the single edge $0\to n$ or visits $r_0$. Indeed, every edge in a nontrivial zero-loss path is an aligned $q$-adic interval of length $d<n$, so $d\le n/q$ and $d\mid r_0$. An interval between consecutive multiples of $d$ cannot contain the multiple $r_0$ strictly inside it.

The proper top checkpoint makes the path nontrivial. If it omits $r_0$, it has a nonzero outer edge, and~\eqref{eq:scaledloss} gives~\eqref{eq:escape}.

If $r_0$ is present and $\beta_{p^br_0}=r_0$, its subset of $H$ is one right coset of its stabilizer $K$ of order $r_0$. This $K$ is normal of index $q$ (Lemma~\ref{lem:normal} below). The whole checkpoint is therefore one right coset of $J=K\times P_b$, normal of prime index $q$ in $G$. Any element outside $J$ acts in a single cycle on the $q$ quotient cosets. Before this checkpoint, a nonempty prefix lies in that one $J$-coset and its translate lies in a disjoint coset. After this checkpoint, if such an element stabilized a prefix, its iterates would force all translates of the full coset into that prefix, hence all of $G$. Thus every proper stabilizer lies in $J$, contradicting $B(\mathbf{x})=G$.

It follows that a present $r_0$ must have $\beta_{p^br_0}<r_0$. Its contribution to $D_{\rm top}$ is at least $p^b>p^{b-1}(p-1)$, so~\eqref{eq:outerbound} again gives~\eqref{eq:escape}. When $r_0=1$, this deficient case is impossible because $\beta\ge1$. No stabilizer persistence between different prefixes is used.

\section{Prime power bases and local accounting}\label{sec:bases}
We prove the complete prime-power dependency, including both axes $b=0$ and $a=0$.

\begin{lemma}\label{lem:normal}
Every maximal proper subgroup of a nontrivial finite $t$-group is normal of index $t$.
\end{lemma}
\begin{proof}
Every proper subgroup $L<Q$ of a finite $t$-group is strictly contained in its normalizer. The class equation gives a nontrivial center $Z(Q)$. If $Z(Q)\not\subseteq L$, a central element outside $L$ belongs to $N_Q(L)\setminus L$. Otherwise induction applies to $L/Z(Q)<Q/Z(Q)$ and its strict normalizer lifts to one for $L$. If $L$ is maximal, $N_Q(L)=Q$. The quotient $Q/L$ is a nontrivial $t$-group with no nontrivial proper subgroup, so it is cyclic of prime order $t$.
\end{proof}

\begin{lemma}[Normal prime index block accounting]\label{lem:block}
Suppose $J\lhd K$ has prime index $t$. If an ordering consists of consecutive right $J$-cosets, then every proper-prefix stabilizer is contained in $J$ and equals the stabilizer of its active local prefix in $J$. Its score is
\begin{equation}
Z_K=\sum_C Z_J(\mathbf{x}|_C)+|K|-|J|.\label{eq:block}
\end{equation}
\end{lemma}
\begin{proof}
An element outside $J$ acts transitively on the $t$ quotient cosets. Inside a block, it cannot fix the unique partially occupied coset. At a proper block boundary, it cannot stabilize the nonempty proper set of full quotient cosets. Thus every proper stabilizer lies in $J$. Elements of $J$ preserve every full and empty coset, so precisely the active local stabilizer remains. Every local contribution appears once, except that the final contribution $|J|-1$ is replaced by $|K|-1$.
\end{proof}

The following identity also applies to interleaved local orders. Suppose $S_j\le J$ for every $j<|K|$. At insertion $j$ into $C_j=Jx_j$, let $L_j^{\rm loc}=\Stab_J(T_j\cap C_j)$. Then $S_j\le L_j^{\rm loc}$ and
\begin{equation}
Z_K=\sum_C Z_J(\mathbf{x}|_C)+|K|-|J|
-\sum_{j<|K|}\bigl(|L_j^{\rm loc}|-|S_j|\bigr).\label{eq:local}
\end{equation}
Each local prefix is counted at its active insertion. The last local and global stabilizers are $J$ and $K$, respectively; every displayed deficit is nonnegative.

If equality forces all these deficits to vanish, the first completion of a $J$-coset makes the entire prefix $J$-invariant. No other coset has yet completed; invariance forces all others to be empty. That coset is the first consecutive block. At each subsequent completion before the last, after the previously completed blocks, the same argument forces exactly the next consecutive block. The final coset fills the remainder. This is the first-completion argument used below.

Now let $|Q|=t^e$. A prefix stabilizer acts freely and its order divides both $j$ and $t^e$. Therefore
\begin{equation}
Z_Q\le\sum_{j=1}^{t^e}\bigl[\gcd(j,t^e)-1\bigr]
=t^e e\left(1-\frac1t\right).\label{eq:basevalue}
\end{equation}
For $Q=1$, the value is $0$ and there is one ordering. For nontrivial $Q$, choose a maximal subgroup $J$, normal of index $t$ by Lemma~\ref{lem:normal}. Optimal local orders exist by induction on $e$, and Lemma~\ref{lem:block} gives attainment of~\eqref{eq:basevalue}.

Conversely, equality in~\eqref{eq:basevalue} saturates every individual prefix cap. At $j=|Q|/t$, the prefix is one right coset of its stabilizer $J$ of that order. The prime-index coset trap from Section~\ref{sec:escape} forces every proper stabilizer into $J$. The local identity~\eqref{eq:local} and inductive value force every local order optimal and every deficit zero. The first-completion argument then supplies consecutive coset blocks. Induction proves exactly the recursive characterization in both directions. At the first full block the stabilizer is $J$, and all proper stabilizers lie in $J$, so $B(\mathbf{x})=J$. Hence the top is unique, and
\[
A(Q)=t!\sum_{J\in\Max(Q)}A(J)^t\quad(Q\ne1),\qquad A(1)=1.
\]
Apply this to $Q=H$ when $b=0$ and to arbitrary $Q=P$ when $a=0$. It proves both axes and their singleton intersection without a general nilpotent theorem.

\section{Attainment for two arbitrary Sylow factors}\label{sec:attainment}
To finish Theorem~\ref{thm:two-sylow}, induct on $a+b$. The prime-power axes and singleton have just been proved. For $b>0$, choose a maximal subgroup $R<P$, normal of index $p$, and put $J=H\times R$. By induction there are orders in $J$ of value $V_{a,b-1}$. Consecutive right $J$-cosets with independent such local orders have, by Lemma~\ref{lem:block}, score
\[
pV_{a,b-1}+np^b-np^{b-1}=V_{a,b}.
\]
This attains the general upper bound. Every maximal subgroup of $H\times P$ is either $H\times R$ with $R$ maximal in $P$, or $M\times P$ with $M$ maximal in $H$, by splitting. The latter direction has index $q$ and the identical accounting
\[
qV_{a-1,b}+np^b-(n/q)p^b=V_{a,b}.
\]
Thus recursive maximal-subgroup orders attain the value in either direction. This is a sufficient construction; necessity for unrestricted $P$ is not inferred.

\section{Cyclic equality and labeled counting}\label{sec:equality}
Here $P=P_b$ is cyclic. Induct on $a+b$, with both axes supplied by Section~\ref{sec:bases}. For $a,b\ge1$ the numerical upper bound is already~\eqref{eq:upper}. Choose $J=H\times P_{b-1}$, normal of index $p$. Its local theorem has smaller $a+b$. Consecutive $J$-cosets with optimal local orders have score
\begin{equation}
pV_{a,b-1}+np^b-np^{b-1}=V_{a,b},\label{eq:attain}
\end{equation}
so the bound is attained.

For any maximizing order, strict escape implies $B(\mathbf{x})<G$. Choose a maximal proper subgroup $J$ containing it. Splitting and the cyclic subgroup chain show that all maximal subgroups are exactly
\begin{equation}
H\times P_{b-1}\quad\text{of index }p,
\qquad M\times P_b\quad(M\in\Max(H))\quad\text{of index }q.\label{eq:maxtypes}
\end{equation}
A split subgroup with both factors proper can enlarge either and is not maximal. With only the $p$-factor proper, maximality requires $P_{b-1}$; with only the $H$-factor proper, it requires a maximal $M$. Every subgroup in~\eqref{eq:maxtypes} is normal of the indicated prime index. The local theorem is the $(a,b-1)$ case, the $(a-1,b)$ case, or a prime-power base. These are all the proper-subgroup inputs required.

All proper-prefix stabilizers lie in $J$, so~\eqref{eq:local} applies. With $t=[G:J]$, the local value satisfies
\begin{equation}
tM(J)+|G|-|J|=V_{a,b}.\label{eq:bothvalues}
\end{equation}
For index $p$ this is~\eqref{eq:attain}; for index $q$ it is
\[
qV_{a-1,b}+np^b-(n/q)p^b=V_{a,b}.
\]
Equality of the global score forces every local order optimal and every local deficit zero. The first-completion argument proves consecutive $J$-coset blocks; local optimality and induction give the entire recursive hierarchy.

Conversely, any recursive ordering uses one of the top subgroups in~\eqref{eq:maxtypes}. Its local orders are optimal by induction and its value is $V_{a,b}$ by Lemma~\ref{lem:block} and~\eqref{eq:bothvalues}. Every proper stabilizer lies in $J$, while at the first full block the stabilizer is exactly $J$. Thus $B(\mathbf{x})=J$, which recovers the top uniquely. This proves equality in both directions without needing to classify equality at every numerical densification step.

For fixed maximal $J$ of index $t$, the distinct right cosets have $t!$ orders, and each independently has $A(J)$ labeled optimal internal orders. Fixing right-translation coordinates gives $t!A(J)^t$ possibilities; changing representatives introduces no additional factor. Unique top recovery makes different $J$ classes disjoint. Summing over~\eqref{eq:maxtypes} proves~\eqref{eq:count}; the prime-power proof supplies~\eqref{eq:qcount}. For $H=1$, the sum is empty, and $A(P_b)=p!A(P_{b-1})^p$. For the singleton, $A(1)=1$ and no top subgroup is asserted. This completes Theorem~\ref{thm:main}.

\section{Rank two escape and equality}\label{sec:ranktwo}
Now let $P=C_p^2$. Theorem~\ref{thm:two-sylow} already gives the value $V_{a,2}$. To prove Theorem~\ref{thm:ranktwo}, first assume $a\ge1$ and $B(\mathbf{x})=G$. Two branches cover the actual $C_2$ checkpoint set. They do not infer a proper full-$P$ checkpoint from $B=G$.

\paragraph{Prime-index trapping.}
We use the general observation already proved in Section~\ref{sec:escape}: if a proper prefix is a single right coset of a normal prime-index subgroup $J$, then all proper-prefix stabilizers lie in $J$. Before that checkpoint a nonempty prefix lies in one coset, which any element outside $J$ moves to a disjoint coset. After it, invariance under an outside element would force every quotient coset, hence the full group. This uses neither a cyclic $P$ nor an abelian ambient group.

\subsection{A proper full projection checkpoint exists}
Suppose $C_2$ has a proper checkpoint. By~\eqref{eq:outerbound}--\eqref{eq:scaledloss}, every nonzero outer edge loses at least $p(p-1)$, and the zero edges are the aligned $q$-power intervals. A nontrivial zero path from $0$ to $n$ must visit $r_0=n/q$: an aligned interval of length $d<n$ has $d\mid r_0$ and cannot cross $r_0$ strictly inside it.

If $r_0$ is absent, one positive edge gives a loss at least $p(p-1)$. If it is present and $\beta_{p^2r_0}=r_0$, its prefix is a single coset of $K\times P$ for $|K|=r_0$, a normal subgroup of prime index $q$. Trapping contradicts $B=G$. Otherwise the retained top deficit is at least $p^2$. The deficient case is impossible when $r_0=1$. Every permitted case gives a loss at least $p(p-1)$, stronger than~\eqref{eq:rankescape}.

\subsection{There is no proper full projection checkpoint}
Suppose $C_2=\{0,N\}$, where $N=np^2$. List all $C_1$ ranks as
\[
0=k_0<k_1<\cdots<k_t=np,\qquad pk_i\in C_1.
\]
A proper checkpoint has a $P$-projection equal to some line of order $p$; these lines may differ. At $j=pk$ the only higher clip is $0<j<N$, since $d\mid n$ implies $p^2d\mid N$. Thus $h_{\C}(pk)=\gcd(k,n)$. At a level-zero point the top clip is redundant beside the level-one clips. The exact cap becomes
\begin{equation}
U_2(\C)=\sum_{i=0}^{t-1}E_p(k_i,k_{i+1})+(N-np).\label{eq:rankextended}
\end{equation}
The correction occurs only at the final endpoint: the edge formula assigns $p\gcd(np,n)-1=pn-1$, whereas the full prefix contributes $N-1$. Throughout this extended interval, the base in $\gcd(k,n)$ and $\D$ remains $n=q^a$, not $np$.

The extended-rank identity~\eqref{eq:loss}--\eqref{eq:zero} gives nonnegative edge losses, exactly zero when $l-k=\gcd(k,l,n)$, and at least $p-1$ otherwise. Periodicity gives $s(np)=p\,s(n)$, and hence
\begin{equation}
F_p(np)+(N-np)=V_{a,2}.\label{eq:rankpotential}
\end{equation}
Every zero edge has aligned length $d\mid n$ and cannot cross rank $n$ strictly inside it. A zero-loss path from $0$ to $np$ must therefore visit the proper rank $n$.

If rank $n$ is absent,~\eqref{eq:rankextended}--\eqref{eq:rankpotential} give a loss at least $p-1$. If it is present and $\beta_{pn}=n$, the actual stabilizer is $H\times L$ for a line $L\le P$. Its order equals the prefix size $pn$, so that prefix is one right coset of $H\times L$, normal of index $p$. Trapping again contradicts $B=G$. Otherwise this endpoint contributes an actual deficit
\[
p(n-\beta_{pn})\ge p
\]
below its pointwise cap. It remains subtractable from $U_2$, giving $Z_G\le V_{a,2}-p$. Both alternatives imply~\eqref{eq:rankescape}. No line-direction switching cost or common subgroup chain has been assumed.

\subsection{Recursive characterization and count}
Induct on $a$, using the finite $p$-group theorem for $a=0$. If $a\ge1$, strict escape forces every optimizer to have $B<G$. A maximal subgroup $J$ containing $B$ is either $H\times L$, for one of the $p+1$ lines $L\le C_p^2$, or $M\times C_p^2$ for $M\in\Max(H)$. These exhaust the possibilities by splitting; their indices are $p$ and $q$, and each is normal. The local structural theorem is respectively the cyclic theorem at $b=1$ or the induction hypothesis at $a-1$.

For $t=[G:J]$, the local value satisfies $tM(J)+|G|-|J|=V_{a,2}$. In the exact local identity~\eqref{eq:local}, every deficit is nonnegative. Global equality forces each local order optimal and all deficits zero. At the first completion of a coset, its local stabilizer is $J$, hence so is the global stabilizer. The prefix is a union of full cosets; all other cosets must be empty because none has completed. Repeating at later completions forces consecutive blocks. Local induction gives precisely the recursive orders.

Conversely those orders attain the maximum by block accounting. Every proper stabilizer lies in the chosen top $J$, and the first full block has stabilizer exactly $J$. Thus $B=J$ recovers the top uniquely. For each fixed $J$ of index $t$ there are $t!A(J)^t$ labeled choices, with no extra representative factor. Distinct top classes are disjoint. Summing over the $p+1$ lines and the maximal $M$ proves~\eqref{eq:rankcount} and finishes Theorem~\ref{thm:ranktwo}.

\section{Why cyclic checkpoint forcing fails in rank two}\label{sec:switch}
Let $D_8=\langle r,s:r^4=s^2=1,\ srs=r^{-1}\rangle$, encode $r^is^j$ by $(i,j)$, and put $P=\mathbb F_3^2$. Set
\[
L=\langle(1,0)\rangle,\quad M=\langle(0,1)\rangle,\quad
A=\{1,s\}\times L,\quad Q=\{r^2,r^3s\}\times M.
\]
The sets $A,Q$ are disjoint and each has six elements. Let $R$ be the six elements of $\{r^2\}\times P$ outside $Q$. Order the group in four sections
\[
A;\quad R;\quad G\setminus(A\cup Q\cup R);\quad Q,
\]
using lexicographic order on $(i,j,u,v)$ within each section, with $0\le i<4$, $0\le j<2$, and $0\le u,v<3$. This specifies a single labeled order without an external list.

At $j=6$, $T_6=A$ and $S_6=\langle s\rangle\times L$. At $j=66$, $T_{66}=G\setminus Q$ and $S_{66}=\langle rs\rangle\times M$, since $Q=(\langle rs\rangle r^2)\times M$. The elements $s,rs$ generate $D_8$, and $L,M$ span $P$, so these proper stabilizers already generate $G$.

There is nevertheless no proper full-$P$ checkpoint. During $A$ a row is partial; during $R$ the initial $A$ rows stay partial. From the end of $R$ through index $66$, the $r^2$ row has exactly six points. During the final $Q$ section, the $r^2$ row completes first and the $r^3s$ row remains partial until the final element. Every nonempty proper prefix therefore has a partial $P$-fibre.

Direct translations in the accepted source give
\begin{equation}
Z=113,\quad V_{3,2}=204,\quad U_2=194,\quad D_{\rm top}=0,
\quad |B|=72,\quad C_2=\{0,72\}.\label{eq:switchvalues}
\end{equation}
This disproves the transferred cyclic implication that $B=G$ forces a proper full-$P$ checkpoint. Its score is well below the maximum; it is neither a counterexample to the present theorems nor a sharp escape witness.

\section{Why nearest level clipping does not suffice}\label{sec:diagnostic}
The use of every enclosing clip is essential to this proof method. For $n=8$, $q=2$, $p=3$, $b=3$, consider the numerical hierarchy
\[
C_1=C_2=C_3=\{0,27,54,\ldots,216\}.
\]
In the first outer interval, through index $27$, a weakened cap testing only $s=r(j)+1$ has value $55$. Densifying $C_1$ lowers that weakened cap to $48$, contradicting the needed monotonic direction. Explicitly, the old open cap is $11\cdot1+5\cdot2+2\cdot4=29$, with endpoint cap $26$. After densification, eight proper level-one nodes give baseline $16$ and $q$-weight increment $6$, for $16+6+26=48$.

Under the actual all-enclosing invariant, the level-three clip suppresses all old open $q$-weights. The old cap is $26$, the first densification gives $16+26=42$, and subsequent elimination reaches $54$. The correct comparison is preserved. This is a counterexample to the weakened proof method, not to the group theorem; no group realizability of this numerical hierarchy is assumed.

\section{The stronger historical dihedral result}\label{sec:d8}
In this section only, let $G=D_8\times P$ and $P=C_p$, so $b=1$. The earlier independently accepted theorem for $D_8\times C_p$, with $p$ odd, proves
\begin{equation}
B(\mathbf{x})=G\quad\Longrightarrow\quad
Z_G(\mathbf{x})\le18p-7=(20p-8)-(2p-1).\label{eq:d8escape}
\end{equation}
Here $V_{3,1}=20p-8$. This specialized loss $2p-1$ is stronger than the general loss $p-1$; the present theorem does not replace or weaken it. Its separate proof uses the complete symbolic 36-edge table on ranks $0,\ldots,8$ and the index-two checkpoint trap. The unchanged theorem and acceptance review are included among the historical source records. Equation~\eqref{eq:d8escape} is reported as that prior result, not derived from~\eqref{eq:escape}.

For its explicit labeled count, put $f=p!$ and
\[
T=f\,2^p+2f^2,\qquad U=f\,8^p+2T^2,\qquad V_4=f\,24^p+6T^2.
\]
These are respectively the counts for $C_2\times C_p$, $C_4\times C_p$, and $C_2^2\times C_p$. Since $D_8$ has one cyclic and two Klein-four maximal subgroups,
\[
A(D_8)=2\cdot8^2+4\cdot24^2=2432,
\]
and the historical recurrence gives
\begin{equation}
A(D_8\times C_p)=p!\,2432^p+2U^2+4V_4^2.\label{eq:d8count}
\end{equation}
At $p=3$, the value is $52$ and the count is $203{,}047{,}403{,}520$.

There is also a uniform escaping example. Write
$D_8=\langle r,s:r^4=s^2=1,\ srs=r^{-1}\rangle$. The order
\[
e,s,r,rs,r^2,r^2s,r^3,r^3s
\]
has proper-prefix stabilizers generating $D_8$ and score $10$. Expanding each element to its full $P$-fibre gives $B(\mathbf{x})=D_8\times P$ and
\[
Z_G=p\cdot10+8(p-1)=18p-8.
\]
This example is one below the specialized upper bound. Neither that observation nor the general theorem asserts the sharp escaping maximum. It also shows why unconditional accumulated-coset containment cannot be used outside its proved hypotheses.

\section{Finite controls separate from the proof}\label{sec:controls}
Sections~\ref{sec:checkpoints}--\ref{sec:ranktwo} give symbolic proofs in exactly the three stated scopes, for all indicated exponents and distinct primes. Finite enumeration is not a premise. The frozen author and independent review packets contain separately written corroborating controls; they were inspected and source-bound for this presentation, not executed during typesetting.

The author packet reports $336$ arbitrary numerical hierarchies through $b=4$, $4{,}325$ whole-gap comparisons, $336$ exact quotient and outer-reduction checks, and $36{,}459$ one-level outer edges with scaled checks through $b=5$. Its $33$ direct group-order checks reach order $216$, including $D_8\times C_{27}$ and the order-$27$ Heisenberg group times $C_8$. Constructed escaping orders cover both the missing-$r_0$ and deficient-$r_0$ branches.

The independent acceptance packet reports $1{,}015$ hierarchies, $1{,}011$ quotient identities, $15{,}582$ whole-gap comparisons including $1{,}650$ gaps longer than $n$, and $155{,}520$ outer edges through depth $b=8$, with no negative-loss failure. It checks $170$ direct orders across $D_8\times C_{27}$, $C_2^3\times C_{27}$, the order-$27$ Heisenberg group times $C_8$, $C_9\times C_{16}$, and $C_{32}$. Three orders have $B=G$, including both escape branches for $D_8\times C_{27}$. The weakened nearest-only diagnostic and the correct all-enclosing comparison are independently reproduced.

An independent dynamic program over every subset uses
\[
D(S)=|\Stab(S)|-1+\max_{x\in S}D(S\setminus\{x\}),\qquad D(\varnothing)=0,
\]
with counts summed over maximizing last elements. It does not assume the optimizer hierarchy. Table~\ref{tab:controls} records the four cyclic controls from the all-depth review. The counts agree with separately calculated recursive-family counts. These cyclic checks establish finite corroborations only; they do not imply any unrestricted quantifier. Separate noncyclic controls are reported below.

\begin{table}[htbp]\centering
\caption{Independent exhaustive cyclic controls in the all-depth review}\label{tab:controls}
\begin{tabular}{lrrr}\toprule
Group & Subset states & Maximum & Maximizing labeled orders\\\midrule
$C_6$ & $64$ & $7$ & $120$\\
$C_{12}$ & $4{,}096$ & $20$ & $31{,}872$\\
$C_{18}$ & $262{,}144$ & $33$ & $13{,}727{,}232$\\
$C_{24}$ & $16{,}777{,}216$ & $52$ & $2{,}044{,}231{,}680$\\\bottomrule
\end{tabular}\end{table}

\paragraph{Noncyclic independent controls.}
The noncyclic acceptance reports $458{,}640$ arithmetic edges, including $423{,}816$ beyond the original rank interval; $245$ arbitrary nested hierarchies; $210$ densification and quotient checks; and $360$ exact extended-rank cap checks. The $524$ base group-order evaluations across $13$ factor specifications include $132$ recursive orders, $6{,}699$ strict clips and $5{,}442$ direct closure/core calculations. Another $30$ evaluations test $D_8$ times the nonabelian order-$27$ Heisenberg group in both orientations; all $14$ recursive examples attain $756$. Four additional escape evaluations cover both alternatives in both rank-two proof branches. One further control has genuinely incomparable proper endpoint lines in $C_2^2\times C_3^2$. The total is $559$ evaluations, including a repeated diagnostic, not $559$ distinct orders.

Table~\ref{tab:rankcontrols} records the separate independent exhaustive subset controls. Their maximizing counts agree with independently calculated recursive counts. No full-$B$ order exists in these four tests, so an absent escape maximum is not a numerical sharpness claim. None is an exhaustive search on $D_8\times C_3^2$.
\begin{table}[htbp]\centering
\caption{Independent exhaustive rank-two controls}\label{tab:rankcontrols}
\begin{tabular}{lrrr}\toprule
Group & Subset states & Maximum & Maximizing labeled orders\\\midrule
$C_2^2$ & $16$ & $4$ & $24$\\
$C_3^2$ & $512$ & $12$ & $5{,}184$\\
$C_3\times C_2^2$ & $4{,}096$ & $20$ & $169{,}344$\\
$C_2\times C_3^2$ & $262{,}144$ & $33$ & $95{,}219{,}712$\\\bottomrule
\end{tabular}\end{table}
For comparison, $D_8\times C_3^2$ has maximum $204$ and seven possible maximal tops in its optimizer family: four of index $3$ and three of index $2$. The order in Section~\ref{sec:switch} has score $113$. These are numerical illustrations of separately proved statements. All reported finite calculations are inherited from the sealed sources; no mathematical source program was run during typesetting.

\section{Related work and the contribution boundary}\label{sec:related}
The mathematical increment beyond version 2 is the closure/core transfer of the value theorem to two arbitrary Sylow factors, together with rank-two escape and optimizer structure. The cyclic-prime-power result remains part of the manuscript. Unique top recovery and the labeled recurrence are consequences of that structure. Independent correctness review and historical comparison answer different questions. The earlier $b=1$ historical comparison and a subsequent bounded all-depth implication test are summarized below. Neither is a priority audit of the new noncyclic results.

For a common labeling, a regular-action permutation $g$ has an invariant proper initial segment exactly at a connectivity-set cut in the terminology of Stanley~\cite{stanley}. Hence $Z_G$ is the sum of connectivity-set sizes over all group elements, including the identity's $|G|-1$ cuts. The present optimization concerns simultaneous relabelings of the entire fixed regular group, preserving multiplication relations. The inspected connectivity statements in the bounded comparison do not supply the earlier $H\times C_p$ maximum or its equality cases.

Fixed-chain transversal signatures are classical; Gonz\'alez Vasco, R\"otteler, and Steinwandt~\cite{gvs} give such constructions and minimum-length signatures for solvable groups. Their length objective is the sum of factor-block sizes. Coset-dependent constructions also predate this work: Baumeister and de Wiljes~\cite{bdw}, Algorithm~5.2 and Theorem~5.4, combine local signatures indexed by transversal choices under stated compatibility conditions. Independent branching or coset-dependent choice is therefore not claimed as an original idea here. Those inspected statements concern factorization blocks, rather than the score over every nested element prefix, and were not found to supply the earlier prefix upper bound or all-optimizer necessity theorem.

Likewise the sum over top choices, factorial ordering of children, and product over independent subproblems are standard counting mechanisms. Maximal-chain enumeration and rooted-tree recursions are relevant antecedents~\cite{st,mbkk}. The objective-specific input to~\eqref{eq:count} is identifying all optimizers and proving disjointness across top subgroups; no new counting principle is claimed.

The bounded $b=1$ comparison located no exact earlier equivalent. That result cannot be promoted to an all-depth or noncyclic priority conclusion. The comparison with Erd\H{o}s, Faigle, and Kern~\cite{efk} remains incomplete: Theorem~2.1 and its main incidence calculation were recovered, but full definitions and Lemma~2.1, most of Section~4, the proof of Proposition~5.1, and the complete Section~6 ending remain partly unread. The checked direct substitutions fail the required hypotheses; this is not a whole-paper exclusion. The preserved historical report records the exact search and access boundaries. No worldwide-first claim, sharp escape result, or general nilpotent theorem follows from this manuscript.

\subsection{The bounded all depth implication test}
The later comparison isolates two failed direct reductions for the cyclic family; it does not exclude every possible classical derivation. Put $\alpha=a(1-1/q)$ and $\eta=b(1-1/p)$. For a cyclic-factor prefix stabilizer $|S_j|=q^{u_j}p^{r_j}$, define $c_0=e_0=1$, $c_i=q^i-q^{i-1}$ and $e_s=p^s-p^{s-1}$ for positive indices, and
\[
n_{i,s}=\#\{j:1\le j\le N,\ i\le u_j,\ s\le r_j\}.
\]
Divisor telescoping and freeness give exactly
\[
Z=\sum_{(i,s)\ne(0,0)}c_ie_s n_{i,s},\qquad
n_{i,s}\le\frac{N}{q^ip^s}.
\]
Independent incidence caps consequently give only $N(\alpha+\eta+\alpha\eta)$. If
\[
\Delta=\sum_{(i,s)\ne(0,0)}c_ie_s
\left[\frac{N}{q^ip^s}-n_{i,s}\right],
\]
then the target upper bound is equivalent to the additional compatibility inequality $\Delta\ge N\alpha\eta$. This inequality is proved here by the checkpoint argument, but the bounded comparison did not derive it from a matched classical incidence theorem. The EFK source gaps stated above therefore remain relevant; no whole-paper exclusion is asserted.

Minimum-length logarithmic signatures also do not by themselves characterize prefix-optimal flattenings. In additive $C_{24}$, the two signatures
\[
([0,8,16],[0,4],[0,2],[0,1]),\qquad
([0,1,2],[0,3],[0,6],[0,12])
\]
have the same ordered type $(3,2,2,2)$ and minimum length $9$, yet last-coordinate-fastest flattening gives scores $23$ and $52$. The first is ordinary mixed-radix order $0,1,\ldots,23$; its proper intervals have trivial translation stabilizer because $I\setminus(I+1)=\{0\}$. The second consecutively traverses a prime-index subgroup chain, so block accounting gives $52$. The common lower length $3+3\cdot2=9$ follows from the inspected minimum-length statement of~\cite{gvs}. This is a counterexample to that specified decoder, not to every generalized signature reduction. The complete bounded comparison is preserved in the package. No classical counting mechanism, independent branching construction, or worldwide-first claim is asserted as new.

\appendix
\section{Version 3 changes and claim boundaries}\label{app:changes}
Version 1 treated $H\times C_p$; version 2 treated $H\times C_{p^b}$ at every cyclic depth. Both source documents, PDFs and packages are preserved byte for byte. This candidate version makes the following source-bound changes.
\begin{enumerate}
\item Theorem~\ref{thm:two-sylow} and Lemma~\ref{lem:closure} add the accepted arbitrary-$P$ maximum and its closure/core bridge. The complete numerical hierarchy argument remains in Sections~\ref{sec:clips}--\ref{sec:edge}; it is not delegated to conversation history or unexplained prior work.
\item Theorem~\ref{thm:main} retains the accepted cyclic equality theorem, sufficient gap and labeled recurrence. The finite prime-power bases and local accounting remain explicit.
\item Theorem~\ref{thm:ranktwo} and Section~\ref{sec:ranktwo} add precisely the accepted rank-two structure. Its two escape branches include the extended-rank edge identity and the retained endpoint deficit. Section~\ref{sec:switch} specifies the accepted direction-switch diagnostic.
\item Sections~\ref{sec:controls} and~\ref{sec:related} distinguish symbolic proofs, inherited finite checks, and bounded historical comparisons. The general-$P$ structure, three-prime extension, sharp escape maxima and historical priority remain outside the claims.
\end{enumerate}
No mathematical expansion beyond the accepted records is intended. Editorial integration and numerical illustrations do not constitute a new independent mathematical acceptance.

\section{Exact source provenance}\label{app:provenance}
The package records every source byte count and SHA-256, together with unchanged copies of both older versions. Archive receipt and Library verification are separate from the payload to avoid self-referential hashes. Mathematical acceptance, document-fidelity review and historical-priority research are distinct statuses.

\paragraph{Accepted noncyclic source.}
The independent acceptance archive is Library version $0$,
\begin{center}\hash{libfile_6b306a62e6348191a2f106e3b19a8e44}\end{center}
It is $233{,}898$ bytes, with SHA-256
\begin{center}\hash{be001a5b8ea0d2e0f546cdb7967f5155e0a2362935d29456a16150f1ac90f76f}\end{center}
It accepts the two-Sylow value, rank-two structural theorem and diagnostic without mathematical repair. Its reviewed author proof is $18{,}250$ bytes, SHA-256
\begin{center}\hash{148ddbc0f0d95ad5d0bde66acd901b4cf2fd69c79283e9c243e7ede0e13d22e0}\end{center}
The author archive is Library version $0$,
\begin{center}\hash{libfile_b790b12b821c8191b50ec98836c8c6c0}\end{center}
with SHA-256
\begin{center}\hash{b1b5b8c2fd41d2a7f2334eba82984abe596bc7f10b8c8cfff1335f3ff00994c4}\end{center}
Its original candidate wording is preserved; the separate review supplies its accepted mathematical status.

\paragraph{Accepted cyclic dependency.}
The all-depth author archive is Library version $0$,
\begin{center}\hash{libfile_1dae9a88f5108191a048f86017d447d6}\end{center}
with SHA-256
\begin{center}\hash{f41867be27ab0f77376b6745adff35575aa3af5044ce4f29d201f60179858d3e}\end{center}
Its independent acceptance is Library version $0$,
\begin{center}\hash{libfile_672d090b20a4819198c511cd125dd489}\end{center}
with SHA-256
\begin{center}\hash{efddb76d4d402aebedf75e1c095ed7c22776ab90142966c0a3a64049cb94f62d}\end{center}
The complete $21{,}879$-byte theorem has SHA-256
\begin{center}\hash{e13674de3c970bc3a18417d9f6e01630b388a0632571637e97674c526f56afd9}\end{center}
This numerical dependency applies to arbitrary admissible hierarchies; its later cyclic-only escape inference is not used for Theorem~\ref{thm:two-sylow}.

\paragraph{Version 2 preserved.}
The independently fidelity-cleared PDF and editable LaTeX remain Library version $0$:
\begin{center}\hash{libfile_6ac21ef7a368819183882483ff5eacb4}\\[3pt]
\hash{libfile_fc51b0ce608c81919a9ec0d80f71422a}\end{center}
Their SHA-256 values are respectively
\begin{center}\hash{aef2c3201b84afb60604299aba4b167e5c14d61e3ee65f0d33ae56cb9579606f}\\[3pt]
\hash{e316918613769298adba4af02ba088ded68062ad961bc505afdd41ac5563ffb8}\end{center}
The original package is
\begin{center}\hash{libfile_61bf9935ca148191b1777fbc24e73447}\end{center}
with SHA-256
\begin{center}\hash{0ba2dcd864e054ee1b89a8672354aba87ee7b438fae98b735da6f3102fbbae3b}\end{center}

\paragraph{Version 1 preserved.}
The PDF, LaTeX and package remain Library version $0$:
\begin{center}\hash{libfile_78c45eb171c08191b2f53963f2aa4845}\\[3pt]
\hash{libfile_ea8ce3f6559481919dfddff2ec859e21}\\[3pt]
\hash{libfile_d2e3566fd3e08191b6c589e14cd43813}\end{center}
Their SHA-256 values are respectively
\begin{center}\hash{24448fdfa5652d6c6c8a182cab38c035cd39d96bba37e1b2abf6f5e65c884b19}\\[3pt]
\hash{cb9cd8f829ce0efc1bec90eed3f5ef896224238515be25fc204024055aacb9d0}\\[3pt]
\hash{245121a1567f6dc4bf35b2b29b775dadae108e968e3bebb1a3e67f0a42cac36d}\end{center}

\paragraph{Historical comparison boundary.}
The bounded $b=1$ comparison is Library
\begin{center}\hash{libfile_07a6103f05bc8191bb3dde2b717f4fc7}\end{center}
The subsequent all-depth known-bound implication comparison is Library
\begin{center}\hash{libfile_5f9f8e18ac9c8191b7719d6428bddf23}\end{center}
with SHA-256
\begin{center}\hash{20d8af204945f115fc5958d5f8991601b6ca3b7513136b55666bc3e331f9a453}\end{center}
These comparisons establish bounded findings only. Their source ledgers and the historical dihedral result remain preserved with the older packages. Third-party full texts are not reproduced in this package.

\begin{thebibliography}{9}
\bibitem{stanley} R. P. Stanley, \emph{The Descent Set and Connectivity Set of a Permutation}. Primary manuscript dated 16 August 2005, equation~(1), Proposition~1.1, and Theorems~1.2 and~1.4. \url{https://math.mit.edu/~rstan/papers/decset.pdf}.
\bibitem{gvs} M. I. Gonz\'alez Vasco, M. R\"otteler, and R. Steinwandt, \emph{On Minimal Length Factorizations of Finite Groups}, Experimental Mathematics 12 (2003). Definition~2.1 and Proposition~3.1. \url{https://www.microsoft.com/en-us/research/wp-content/uploads/2016/11/On-Minimal-Length-Factorizations-of-Finite-Groups.pdf}.
\bibitem{bdw} B. Baumeister and J.-H. de Wiljes, \emph{Aperiodic logarithmic signatures}, arXiv:1104.4503 (2011), Algorithm~5.2 and Theorem~5.4. \url{https://arxiv.org/abs/1104.4503}.
\bibitem{st} M. \c{S}tef\u{a}nescu and M. T\u{a}rn\u{a}uceanu, \emph{Counting maximal chains of subgroups of finite nilpotent groups}, Carpathian Journal of Mathematics 25 (2009), inspected pages 119--124, Theorem~2.1 and Section~3. \url{https://www.carpathian.cunbm.utcluj.ro/wp-content/uploads/2009-vol-25-1/carpathian_2009_25_1_119_127.pdf}.
\bibitem{mbkk} F. A. Matsen IV, S. C. Billey, A. Kas, and M. Konvalinka, \emph{Tanglegrams A Reduction Tool for Mathematical Phylogenetics}, primary manuscript, Section~2.2, Theorem~5 and Example~6. \url{https://matsen.fredhutch.org/papers/Matsen2018-bg.pdf}.
\bibitem{efk} P. Erd\H{o}s, U. Faigle, and W. Kern, \emph{A Group-Theoretic Setting for Some Intersecting Sperner Families}, Combinatorics, Probability and Computing 1 (1992), 323--334. Partially inspected; the remaining access boundary is stated in Section~\ref{sec:related}. \url{https://doi.org/10.1017/S0963548300000377}.
\end{thebibliography}
\end{document}
